Redistricting ensembles are used to test whether an enacted map is
a statistical outlier.
But how large must an ensemble be for the test to have adequate power?
This question is rarely asked; experts routinely run 10,000 or 50,000
steps without demonstrating that the run is sufficient to detect the
size of gerrymander alleged.

We derive a power function for redistricting tests as a function of
ensemble size, chain autocorrelation, null distribution variance,
and effect size (partisan advantage in percentage points).
The key input is the G.4~\citep{g4paper} ESS formula: for a
\recom{} chain with lag-1 Hamming autocorrelation $\rho_1$,
a chain of $n$ steps yields $\ess \approx n(1 - \rho_1)/(1 + \rho_1)$.

For North Carolina ($k = 14$, $\rho_1 = 0.87$ from G.4):
80\% power to detect a 10-percentage-point partisan advantage at
$\alpha = 0.05$ requires $\ess \approx 385$ independent samples,
corresponding to approximately $n \approx 5{,}000$ chain steps.

For California ($k = 52$, $\rho_1 = 0.94$):
the same power requires $\ess \approx 385$, corresponding to
$n \approx 30{,}000$ chain steps due to the higher autocorrelation.

Current practice of 10,000 steps gives 80\% power for NC (ESS = 695)
but only 45\% power for California (ESS = 588 at 10,000 steps, but
$\rhat = 1.09$ indicating the chain has not yet mixed---so the ESS
is not meaningfully computed).
We provide a 50-state minimum ensemble table and a litigation-grade
recommendation: 20,000 steps for states with $k \leq 17$ districts,
50,000 steps for $k > 17$.
At \bisect{}'s H.2 throughput of 50,000 steps per second for the
ensemble module, these are feasible in under 2 seconds for any state.
